\scanpage{531}
\mspage{442}
% archive sheet number: 530

Notons tout de suite deux corollaires

\uline{Corollaire} 4.15.7.2. Supposons que $C$ soit sans projecteurs. Alors $\underline{x}=(x_i)_{i\in \unsure{J|I}}$ est un ind-objet ess. constant sss il existe un foncteur cofinal $J\longrightarrow I$, avec $J$ \deleted{\unsure{ordonné}} \inserted{filtrant}, de telle façon que le système inductif réindexé $(x_j)_{j\in J}$ soit \deleted{isomorphe (\uline{dans la catégorie} $\underline{\mathrm{Hom}}(J,C)$) à un système inductif} \uline{constant} (\unsure{défini} par un foncteur constant $J\longrightarrow C$ ; tous les $x_j$ égaux à un $=x$, et les morphismes de \deleted{\unsure{trans|tous}} tous des identités).
% A broken line runs under the crossed wording from isomorphe to inductif.

% slanted
\margin{\unsure{et OPS|cf. OPS|cf. OP 3} $J=i_0\backslash I\xrightarrow[\mathrm{can.}]{} I,$}

La suffisance est claire, pour la nécessité on applique 4.15.7.1, qui nous dit qu'on peut réindexer de façon à satisfaire aux conditions $1^\circ)$ à $4^\circ)$ (\deleted{pour} des projecteurs $p_j$ convenables \unsure{des|de les} $x_j$). Mais l'hyp. sur $C$ implique que \deleted{les} $p_j \unsure{=|\Rightarrow}\mathrm{id}_{x_j}\quad\forall j\in\mathrm{Ob}\,J$, et les conditions $1^\circ,2^\circ,3^\circ$ sont satisfaites tautologiquement pour un tel choix, la condition $4^\circ$ signifie que les morphismes de transition sont des isomorphismes.

\deleted{\unsure{OPS|OP 3} que $J$ soit ordonné. Mais il est immédiat \deleted{pour} -- cela et pour un système inductif $(x_j)_{j\in J}$ donné, avec $J$ \uline{1-connexe}, \inserted{(\unsure{p. ex.|cf. car} $J$ \unsure{filtrant|filtr.})}, la condition que les $x_\lambda$ soient des iso est néc. et suff. pour que ce syst. soit isom. à un syst. ind. constant, d'où la conclusion}

% slanted
\margin{\unsure{Garder les traits biffés|Garder le texte biffé}}
